% !TeX root = ../sections/02_exercises.tex
\subsection{2-A-2.}
\begin{exercise}
	関数 \(f:X\to\mathbb{R}\) が有界であることは，
	\[
	\exists M>0\quad \text{s.t.}\quad
	\forall x\in X,\quad |f(x)|\leq M
	\]
	であることと同値であることを証明せよ。
\end{exercise}

\begin{background}
	関数 \(f:X\to\mathbb{R}\) が有界であるとは，
	値域
	\[
	f(X)=\{f(x)\mid x\in X\}
	\]
	が \(\mathbb{R}\) の部分集合として有界であることをいう。
	
	すなわち，ある実数 \(m,M\) が存在して，
	任意の \(x\in X\) に対して
	\[
	m\leq f(x)\leq M
	\]
	が成り立つことである。
	
	一方，
	\[
	|f(x)|\leq C
	\]
	は
	\[
	-C\leq f(x)\leq C
	\]
	と同値である。
	したがって，絶対値による有界性は，
	上下両側からの有界性を一つの式で表している。
\end{background}

\begin{strategy}
	まず，
	\[
	\exists M>0,\quad |f(x)|\leq M
	\quad (x\in X)
	\]
	を仮定する。
	このとき
	\[
	-M\leq f(x)\leq M
	\]
	であるから，\(f(X)\) は上にも下にも有界である。
	したがって \(f\) は有界である。
	
	逆に，\(f\) が有界であると仮定する。
	すると，ある実数 \(m,M\) が存在して
	\[
	m\leq f(x)\leq M
	\qquad (x\in X)
	\]
	が成り立つ。
	
	ここで
	\[
	C=\max\{|m|,|M|,1\}
	\]
	とおく。
	すると \(C>0\) であり，任意の \(x\in X\) に対して
	\[
	|f(x)|\leq C
	\]
	が成り立つ。
	
	よって，二つの条件は同値である。
\end{strategy}